%% kapa: exam=1 round=33 year=2022 period=1 form=A subject=economics subject_ko=경제학원론 seq=12 number=52
%% source: Q-Net 감정평가사 기출문제 (cst00309) artlSeq=5209997 file=raw/1008/5209997/2022년도 제33회 감정평가사 1차시험 1교시(민법,경제학,부동산학).pdf
%% conversion: kordoc 4.24.51 -> tools/gen1.py (auto); review=proofread-sonnet; flags=none
\begin{question}{52}
( )에 들어갈 내용으로 옳은 것은? (단, P는 가격, Q는 수요량이다.)
\begin{framedtext}
독점기업의 수요곡선은 P = 30 $-$ 2Q이고 현재 가격이 10이다. 이 때 수요의 가격탄력성은 ( ㄱ )이고, 총수입을 증대시키기 위해 가격을 ( ㄴ )해야 한다.
\end{framedtext}
\begin{choices}
  \item ㄱ: 비탄력적, ㄴ: 인하
  \item ㄱ: 비탄력적, ㄴ: 인상
  \item ㄱ: 단위탄력적, ㄴ: 유지
  \item ㄱ: 탄력적, ㄴ: 인하
  \item ㄱ: 탄력적, ㄴ: 인상
\end{choices}
\end{question}
\answer{2}
